% appendices/detailed_results.tex
% Apéndice C: Resultados detallados de pruebas
% ============================================================================
%
% GUÍA: Incluir aquí los datos extensos que no caben en el cuerpo (40 pp):
%
%   \section{Resultados de pruebas unitarias}
%   - Listado completo de tests ejecutados (nombre, resultado, tiempo).
%   - Informe de cobertura por fichero (salida de coverage.py).
%   - Incluir como código o tabla larga:
%     \Code[cod:coverage_report]{Informe de cobertura}
%          {Informe de cobertura generado por \texttt{pytest-cov}}
%          {codes/coverage_report.txt}{1}{50}{1}{text}
%     TODO: Exportar el informe a codes/coverage_report.txt
%
%   \section{Resultados de pruebas de integración}
%   - Tabla detallada: endpoint, método, test, resultado, tiempo.
%   - Listado completo de peticiones probadas.
%
%   \section{Resultados de pruebas de carga}
%   - Gráficas de rendimiento (tiempos de respuesta, throughput):
%     \begin{figure}{fig:load_latency}{Latencia bajo carga: distribución de tiempos de respuesta}
%       \image{}{}{load_test_latency}
%     \end{figure}
%     TODO: Exportar gráficas de la herramienta de carga a figures/.
%   - Gráficas de uso de recursos (CPU, memoria, red) si se midieron.
%   - Logs o extractos relevantes.
%
%   \section{Pruebas de aceptación}
%   - Tabla de requisitos vs. resultado de verificación:
%     \begin{table}{tab:acceptance}{Verificación de requisitos funcionales}
%       \begin{tabular}{l l l}
%         \textbf{Requisito} & \textbf{Descripción} & \textbf{Estado} \\
%         \hline
%         RF-1  & Autenticación y sesiones   & Cumplido / Parcial / No cumplido \\
%         RF-2  & Gestión de usuarios        & ... \\
%         ...
%       \end{tabular}
%     \end{table}
% ============================================================================

% TODO: Completar con los resultados reales cuando se ejecuten las pruebas.

[Resultados detallados de pruebas pendientes.]
